% =============================================================
%  TỔNG HỢP CÁC MÔ HÌNH HÀNG ĐỢI – Cơ sở truyền số liệu (ET4070)
%  Biên dịch:  xelatex Tong_hop_Mo_hinh_Hang_doi.tex  (chạy 2 lần)
% =============================================================
\documentclass[11pt,a4paper]{article}

% ---------- Trang & phông chữ ----------
\usepackage[a4paper,margin=2.1cm,top=2.2cm,bottom=2.3cm]{geometry}
\usepackage{fontspec}
\usepackage{amsmath,amssymb}
\usepackage{unicode-math}
\setmainfont{Cambria}
\setsansfont{Calibri}
\setmonofont{Consolas}
\setmathfont{Cambria Math}
\usepackage[vietnamese]{babel}

% ---------- Gói tiện ích ----------
\usepackage{xcolor}
\usepackage{graphicx}
\usepackage{booktabs,tabularx,array,multirow}
\usepackage{enumitem}
\usepackage{needspace}
\usepackage[most]{tcolorbox}
\usepackage{tikz}
\usetikzlibrary{arrows.meta,positioning,calc}
\usepackage{titlesec}
\usepackage{fancyhdr}
\usepackage[hidelinks]{hyperref}
\usepackage{bookmark}

% ---------- Màu sắc ----------
\definecolor{ink}{HTML}{1F2933}
\definecolor{accent}{HTML}{0B5D7A}      % xanh đậm: tiêu đề, khung
\definecolor{accentsoft}{HTML}{E6F1F5}
\definecolor{warm}{HTML}{B4541A}        % cam: cảnh báo, lỗi hay gặp
\definecolor{warmsoft}{HTML}{FBEEE4}
\definecolor{good}{HTML}{2E7D4F}        % xanh lá: ví dụ, lời giải
\definecolor{goodsoft}{HTML}{E8F3EC}
\definecolor{rule}{HTML}{C9D3DA}
\color{ink}

% ---------- Tiêu đề mục ----------
\titleformat{\section}{\sffamily\LARGE\bfseries\color{accent}}{\thesection}{0.6em}{}[\vspace{-0.4em}\color{rule}\rule{\linewidth}{0.8pt}]
\titleformat{\subsection}{\sffamily\Large\bfseries\color{accent}}{\thesubsection}{0.6em}{}
\titleformat{\subsubsection}{\sffamily\large\bfseries\color{ink}}{\thesubsubsection}{0.5em}{}
\titlespacing*{\section}{0pt}{1.2em}{0.9em}
\titlespacing*{\subsection}{0pt}{1.1em}{0.5em}

\setlength{\parindent}{0pt}
\setlength{\parskip}{0.45em}
\setlist{itemsep=0.15em,topsep=0.3em}
\renewcommand{\arraystretch}{1.25}

% ---------- Header / footer ----------
\setlength{\headheight}{14pt}
\pagestyle{fancy}
\fancyhf{}
\fancyhead[L]{\sffamily\small\color{accent}Tổng hợp mô hình hàng đợi}
\fancyhead[R]{\sffamily\small\color{accent!70}Cơ sở truyền số liệu}
\fancyfoot[C]{\sffamily\small\thepage}
\renewcommand{\headrulewidth}{0.4pt}
\renewcommand{\headrule}{\color{rule}\hrule width\headwidth height\headrulewidth}

% ---------- Hộp nội dung ----------
\tcbset{common/.style={enhanced,breakable,boxrule=0.6pt,arc=2pt,left=7pt,right=7pt,top=5pt,bottom=5pt,
  fonttitle=\sffamily\bfseries}}
% Công thức quan trọng
\newtcolorbox{congthuc}[1][]{common,colback=accentsoft,colframe=accent,title={#1}}
% Ví dụ có lời giải
\newtcolorbox{vidu}[1][]{common,colback=white,colframe=good,coltitle=white,colbacktitle=good,title={Ví dụ: #1}}
% Lưu ý / bẫy
\newtcolorbox{luuy}[1][Lưu ý]{common,colback=warmsoft,colframe=warm,title={#1}}
% Nhận diện mô hình
\newtcolorbox{nhandien}{common,colback=white,colframe=rule,borderline west={3pt}{0pt}{accent},
  title={Nhận diện trong đề},coltitle=accent,colbacktitle=white}
% Hộp tóm tắt chung
\newtcolorbox{tomtat}[1][]{common,colback=goodsoft,colframe=good,title={#1}}

\newcommand{\loigiai}{\par\smallskip{\sffamily\bfseries\color{good}Lời giải.}\ }
\newcommand{\dapan}[1]{\fcolorbox{good}{goodsoft}{\ensuremath{#1}}}

% ---------- Sơ đồ chuỗi sinh–tử ----------
% Dùng: \bdchain{0/-/-/st, 1/\lambda/\mu/st, ...}
%   mỗi phần tử = tên trạng thái / nhãn mũi tên đến (từ trạng thái trước) / nhãn mũi tên đi (về trạng thái trước) / kiểu nút
\tikzset{
  st/.style={circle,draw=accent,very thick,fill=accentsoft,minimum size=9.5mm,inner sep=0pt,font=\small},
  dots/.style={draw=none,fill=none,minimum size=9.5mm,font=\small},
  bdarr/.style={-{Stealth[length=2.2mm]},thick,color=ink!80},
  bdlab/.style={font=\footnotesize,inner sep=1.5pt}
}
\newcommand{\bdchain}[2][16mm]{%
\begin{center}
\begin{tikzpicture}[x=#1]
\foreach \lab/\up/\down/\sty [count=\k from 0] in {#2} {
  \node[\sty] (s\k) at (\k,0) {$\lab$};
  \ifnum\k>0
    \pgfmathtruncatemacro{\p}{\k-1}
    \draw[bdarr] (s\p) to[bend left=38] node[above,bdlab]{$\up$} (s\k);
    \draw[bdarr] (s\k) to[bend left=38] node[below,bdlab]{$\down$} (s\p);
  \fi
}
\end{tikzpicture}
\end{center}}

% ---------- Siêu dữ liệu PDF ----------
\hypersetup{pdftitle={Tổng hợp các mô hình hàng đợi},pdfsubject={Cơ sở truyền số liệu – Lý thuyết hàng đợi}}

\begin{document}

% ================= TRANG BÌA =================
\begin{titlepage}
\thispagestyle{empty}
\vspace*{3.2cm}
{\sffamily\color{accent!80}\large CƠ SỞ TRUYỀN SỐ LIỆU · ET4070}\par\vspace{0.6em}
{\color{accent}\rule{0.35\linewidth}{2pt}}\par\vspace{1.4em}
{\sffamily\bfseries\fontsize{34}{40}\selectfont\color{ink} Tổng hợp\\[0.15em] các mô hình hàng đợi\par}
\vspace{1.4em}
{\large Tham số, công thức, cách nhận diện và ví dụ có lời giải\\
cho M/M/1, M/M/1/K, M/M/c, M/M/c/c, M/M/c/K, M/M/$\infty$, M/G/1\\
và mạng hàng đợi (Jackson, Burke, Kleinrock)\par}
\vfill
\begin{tcolorbox}[common,colback=accentsoft,colframe=accent,width=\linewidth]
\textbf{Câu hỏi mở đầu:} \emph{Mỗi mô hình có tham số và công thức khác nhau phải không?}\par
\textbf{Đúng.} Mỗi mô hình khác nhau ở \textbf{số server $c$}, \textbf{sức chứa $K$}, cách định nghĩa \textbf{$\rho$} và \textbf{công thức $p_0$}. Tuy vậy, tất cả đều đi theo \textbf{một khung 5 bước chung} (mục~\ref{sec:khung}). Nắm khung này thì gặp mô hình nào cũng làm được.
\end{tcolorbox}
\vspace{1cm}
{\small\color{ink!70} Ký hiệu bám theo slide bài giảng: $p_0, p_i, \rho, P_Q, \lambda_a, \theta_i, \gamma$.\\
Các số liệu trong ví dụ đã được kiểm tra lại bằng máy tính.}
\end{titlepage}

\tableofcontents
\newpage

% =============================================================
\section{Bức tranh tổng quan}
% =============================================================

\subsection{Các mô hình khác nhau ở đâu?}

Mọi mô hình trong môn học đều mô tả cùng một thứ: \textbf{khách hàng} (gói tin, cuộc gọi, tác vụ...) đi vào, \textbf{xếp hàng} nếu các server đều bận, được \textbf{phục vụ} rồi đi ra. Các mô hình chỉ khác nhau ở bốn điểm:

\begin{center}
\begin{tabularx}{\linewidth}{@{}l X X@{}}
\toprule
\textbf{Điểm khác} & \textbf{Ý nghĩa} & \textbf{Ảnh hưởng tới công thức} \\
\midrule
Số server $c$ & Bao nhiêu khách được phục vụ \emph{cùng lúc} & Tốc độ đi của trạng thái $i$ là $\min(i,c)\,\mu$ thay vì luôn bằng $\mu$ \\
Sức chứa $K$ & Hệ thống chứa tối đa bao nhiêu khách & Có $K$ hữu hạn thì có \textbf{xác suất từ chối} $p_K$ và phải dùng \textbf{$\lambda_a$} \\
Định nghĩa $\rho$ & $\rho = \lambda/\mu$ (1 server) hay $\rho=\lambda/(c\mu)$ ($c$ server) & Đổi điều kiện ổn định và mọi công thức phía sau \\
Phân bố phục vụ & Mũ (M), hằng số (D), bất kỳ (G) & M/G/1 không dùng được chuỗi Markov, phải dùng công thức P-K \\
\bottomrule
\end{tabularx}
\end{center}

\subsection{Bảng ký hiệu dùng xuyên suốt}

\begin{center}
\begin{tabularx}{\linewidth}{@{}c X c@{}}
\toprule
\textbf{Ký hiệu} & \textbf{Ý nghĩa} & \textbf{Đơn vị} \\
\midrule
$\lambda$ & Tốc độ đến trung bình (của luồng Poisson) & khách/đơn vị t.gian \\
$\mu$ & Tốc độ phục vụ của \textbf{một} server; $1/\mu = T_S$ là thời gian phục vụ trung bình & khách/đơn vị t.gian \\
$c$ & Số server (số quầy, số nhân, số kênh) & -- \\
$K$ & Sức chứa \textbf{toàn hệ thống} (gồm cả khách đang được phục vụ) & khách \\
$a = \lambda/\mu$ & Cường độ tải (đơn vị \emph{Erlang}); slide viết là $c\rho$ & Erlang \\
$\rho$ & Hệ số sử dụng của một server: $\rho=\lambda/\mu$ (1 server), $\rho = \lambda/(c\mu)$ ($c$ server) & -- \\
$p_i$ & Xác suất trong hệ thống có đúng $i$ khách (ở trạng thái dừng) & -- \\
$p_0$ & Xác suất hệ thống rỗng & -- \\
$N,\ N_Q$ & Số khách trung bình trong \textbf{hệ thống}, trong \textbf{hàng đợi} & khách \\
$T,\ T_Q$ & Thời gian trung bình ở trong \textbf{hệ thống} (trễ), ở trong \textbf{hàng đợi} (chờ) & đơn vị t.gian \\
$\lambda_a$ & Tốc độ đến \emph{thực tế} (đã trừ khách bị từ chối): $\lambda_a = \lambda(1-p_K)$ & khách/đơn vị t.gian \\
$P_Q$ & Xác suất khách đến phải chờ (công thức Erlang C) & -- \\
$P_B$ & Xác suất khách bị chặn/từ chối (công thức Erlang B) & -- \\
\bottomrule
\end{tabularx}
\end{center}

\begin{luuy}[Hệ thống và hàng đợi khác nhau]
\textbf{Hệ thống} = hàng đợi + các server, nên luôn có $T = T_Q + \dfrac{1}{\mu}$ và $N = N_Q + (\text{số server bận trung bình})$.\\
Khi đề hỏi \emph{“thời gian chờ”} thì thường là $T_Q$. Khi đề hỏi \emph{“trễ”} hoặc \emph{“từ lúc vào đến lúc xong”} thì là $T$.
\end{luuy}

\subsection{Khung 5 bước dùng cho mọi mô hình}\label{sec:khung}

\begin{tomtat}[Quy trình giải một bài hàng đợi]
\begin{enumerate}[leftmargin=1.6em,label=\textbf{B\arabic*.}]
  \item \textbf{Nhận diện mô hình} A/S/$c$/$K$ (xem bảng ở mục~\ref{sec:nhandien}). \textbf{Đổi đơn vị} về cùng một đơn vị thời gian.
  \item \textbf{Tính tham số}: $\lambda$, $\mu$ (nếu đề cho dung lượng kênh $C$ và độ dài gói $L$ thì $\mu = C/L$), $a=\lambda/\mu$, $\rho$. Kiểm tra \textbf{điều kiện ổn định} nếu $K=\infty$.
  \item \textbf{Vẽ chuỗi sinh--tử}, viết $p_i$ theo $p_0$ (cân bằng luồng giữa hai trạng thái kề nhau).
  \item \textbf{Tính $p_0$} từ điều kiện $\sum p_i = 1$. Từ đó suy ra các xác suất đề hỏi: $p_K$, $P_Q$, $P_B$...
  \item \textbf{Tính $N$ hoặc $N_Q$}, sau đó dùng \textbf{định lý Little} để ra $T$, $T_Q$. Nếu có $K$ hữu hạn thì dùng $\lambda_a$ thay cho $\lambda$.
\end{enumerate}
\end{tomtat}

\subsection{Bản đồ quan hệ giữa các mô hình}

Mô hình tổng quát nhất trong môn là \textbf{M/M/$c$/$K$}. Các mô hình còn lại là trường hợp riêng của nó:

\begin{center}
\begin{tikzpicture}[
  box/.style={draw=accent,very thick,fill=accentsoft,rounded corners=2pt,align=center,font=\sffamily\small,minimum width=27mm,minimum height=10mm},
  big/.style={box,fill=accent,text=white,font=\sffamily\bfseries},
  lk/.style={-{Stealth[length=2.2mm]},thick,color=ink!70},
  lb/.style={font=\footnotesize\itshape,fill=white,inner sep=1.5pt}]
  \node[big] (gen) at (0,0) {M/M/$c$/$K$};
  \node[box] (mm1k) at (-5.2,-2.2) {M/M/1/$K$};
  \node[box] (mmcc) at (0,-2.2) {M/M/$c$/$c$\\[-1pt]\footnotesize(Erlang B)};
  \node[box] (mmc)  at (5.2,-2.2) {M/M/$c$\\[-1pt]\footnotesize(Erlang C)};
  \node[box] (mm1)  at (0,-4.6) {M/M/1};
  \node[box] (mminf) at (5.2,-4.6) {M/M/$\infty$};
  \draw[lk] (gen) -- node[lb]{$c=1$} (mm1k);
  \draw[lk] (gen) -- node[lb]{$K=c$} (mmcc);
  \draw[lk] (gen) -- node[lb]{$K\to\infty$} (mmc);
  \draw[lk] (mm1k) -- node[lb]{$K\to\infty$} (mm1);
  \draw[lk] (mmc) -- node[lb]{$c=1$} (mm1);
  \draw[lk] (mmc) -- node[lb]{$c\to\infty$} (mminf);
\end{tikzpicture}
\end{center}

Nhờ quan hệ này, bạn có thể \textbf{tự kiểm tra} công thức. Ví dụ: đặt $c=1$ vào công thức M/M/$c$ phải ra đúng công thức M/M/1.

% =============================================================
\section{Kiến thức nền tảng}
% =============================================================

\subsection{Tiến trình Poisson và phân bố mũ}

Chữ \textbf{M} (Markov, \emph{không nhớ}) trong ký hiệu Kendall đều dựa trên hai sự kiện sau.

\begin{congthuc}[Định lý 1: số lần đến trong khoảng $t$ có phân bố Poisson]
\[
P\{N(t)=k\} = \frac{(\lambda t)^k}{k!}\,e^{-\lambda t},\qquad k=0,1,2,\dots \qquad E[N(t)] = \lambda t
\]
\end{congthuc}

\begin{congthuc}[Định lý 2: khoảng thời gian giữa hai lần đến có phân bố mũ]
\[
P\{X \le t\} = 1-e^{-\lambda t},\qquad P\{X > t\} = e^{-\lambda t},\qquad E[X]=\frac{1}{\lambda}
\]
\end{congthuc}

\textbf{Các tính chất hay dùng:}
\begin{itemize}
  \item \textbf{Không nhớ:} $P\{X > s+t \mid X > s\} = P\{X>t\}$. Đã chờ bao lâu thì thời gian \emph{còn lại} vẫn có cùng phân bố mũ như lúc đầu.
  \item \textbf{Ghép luồng:} gộp các luồng Poisson độc lập $\lambda_1,\lambda_2,\dots$ được một luồng Poisson với tốc độ $\sum\lambda_i$.
  \item \textbf{Rẽ luồng:} mỗi khách được đưa sang nhánh $j$ với xác suất $p_j$ (độc lập nhau) thì nhánh $j$ là luồng Poisson với tốc độ $p_j\lambda$.
\end{itemize}

\begin{vidu}[xe buýt (slide Tiến trình Poisson)]
Xe buýt đến bến theo tiến trình Poisson, trung bình \textbf{5 phút/chuyến}. Một người vừa nhỡ một chuyến. Tính xác suất người đó đợi chuyến tiếp theo \textbf{ít hơn 3 phút}.
\loigiai “5 phút/chuyến” là khoảng cách \emph{trung bình} giữa hai chuyến, nên $\lambda = 1/5 = 0{,}2$ chuyến/phút. Nhờ tính không nhớ, “vừa nhỡ xe” nghĩa là bắt đầu chờ lại từ đầu:
\[
P\{X<3\} = 1-e^{-0{,}2\cdot 3} = 1-e^{-0{,}6} \approx \dapan{0{,}451}
\]
\end{vidu}

\subsection{Định lý Little}

\begin{congthuc}[Định lý Little: đúng cho mọi hệ thống ổn định]
\[
N = \lambda\,T \qquad\qquad N_Q = \lambda\,T_Q \qquad\qquad N_S = \lambda\cdot\frac{1}{\mu}
\]
Nếu hệ thống \textbf{từ chối khách} ($K$ hữu hạn) thì thay $\lambda$ bằng \textbf{$\lambda_a=\lambda(1-p_K)$}:\quad $N=\lambda_a T$, $N_Q = \lambda_a T_Q$.
\end{congthuc}

Định lý \emph{không} phụ thuộc vào phân bố đến, phân bố phục vụ, số server hay quy tắc phục vụ. Biết 2 trong 3 đại lượng (số lượng, tốc độ, thời gian) là suy ra được đại lượng còn lại.

\subsection{Ký hiệu Kendall}

Slide viết đầy đủ là \textbf{A/S/m/[B]/[K]/[SD]}. Tài liệu này dùng dạng rút gọn \textbf{A/S/$c$/$K$}, trong đó $K$ là \emph{sức chứa hệ thống}.

\begin{center}
\begin{tabularx}{\linewidth}{@{}c l X@{}}
\toprule
\textbf{Vị trí} & \textbf{Ý nghĩa} & \textbf{Giá trị thường gặp} \\
\midrule
A & Tiến trình đến & M (Poisson/mũ), D (tất định), $E_r$ (Erlang bậc $r$), $H_r$ (siêu mũ), G (tổng quát) \\
S & Tiến trình phục vụ & M, D, $E_r$, $H_r$, G (như trên) \\
$c$ & Số server & $1, 2, \dots, \infty$ \\
$K$ & Sức chứa hệ thống & Mặc định $\infty$. Có $c$ server và $B$ chỗ chờ thì $K = B + c$ \\
-- & Kích thước quần thể & Mặc định $\infty$ (nguồn khách vô hạn) \\
SD & Quy tắc phục vụ & Mặc định FCFS (FIFO). Ngoài ra có LCFS, SIRO, RR, PS... \\
\bottomrule
\end{tabularx}
\end{center}

\needspace{14\baselineskip}
\subsection{Quá trình sinh--tử: gốc của mọi công thức M/M}

Trạng thái của hệ thống là \textbf{số khách $i$}. Một khách đến (\emph{sinh}) đưa trạng thái từ $i$ lên $i+1$ với tốc độ $\lambda_i$. Một khách đi (\emph{tử}) đưa trạng thái từ $i$ về $i-1$ với tốc độ $\mu_i$.

\bdchain{0/-/-/st, 1/\lambda_0/\mu_1/st, 2/\lambda_1/\mu_2/st, \cdots/\lambda_2/\mu_3/dots, i/\lambda_{i-1}/\mu_i/st, i{+}1/\lambda_i/\mu_{i+1}/st, \cdots/\lambda_{i+1}/\mu_{i+2}/dots}

\begin{congthuc}[Phương trình cân bằng và nghiệm tổng quát]
Ở trạng thái dừng, luồng đi qua giữa hai trạng thái kề nhau phải cân bằng:
\[
\lambda_{i}\,p_{i} = \mu_{i+1}\,p_{i+1}
\quad\Longrightarrow\quad
p_i = p_0\prod_{k=0}^{i-1}\frac{\lambda_k}{\mu_{k+1}},
\qquad
p_0 = \Bigl[\,1+\sum_{i\ge 1}\prod_{k=0}^{i-1}\frac{\lambda_k}{\mu_{k+1}}\Bigr]^{-1}
\]
\end{congthuc}

\textbf{Mọi mô hình M/M/\dots{} chỉ khác nhau ở cách chọn $\lambda_i$, $\mu_i$:}
\begin{center}
\begin{tabular}{@{}lll@{}}
\toprule
\textbf{Mô hình} & \textbf{$\lambda_i$ (tốc độ đến ở trạng thái $i$)} & \textbf{$\mu_i$ (tốc độ đi ở trạng thái $i$)} \\
\midrule
M/M/1 & $\lambda$ với mọi $i$ & $\mu$ \\
M/M/1/$K$ & $\lambda$ nếu $i<K$, bằng $0$ nếu $i=K$ & $\mu$ \\
M/M/$c$ & $\lambda$ & $\min(i,c)\,\mu$ \\
M/M/$c$/$c$ & $\lambda$ nếu $i<c$, bằng $0$ nếu $i=c$ & $i\,\mu$ \\
M/M/$c$/$K$ & $\lambda$ nếu $i<K$, bằng $0$ nếu $i=K$ & $\min(i,c)\,\mu$ \\
M/M/$\infty$ & $\lambda$ & $i\,\mu$ \\
\bottomrule
\end{tabular}
\end{center}

% =============================================================
\section{Các mô hình một server}
% =============================================================

% -------------------------------------------------------------
\subsection{M/M/1 (một server, hàng đợi vô hạn)}
% -------------------------------------------------------------

\begin{nhandien}
Khách đến theo Poisson, thời gian phục vụ có phân bố mũ, \textbf{một} quầy/đường truyền/CPU, chỗ chờ “đủ lớn” hoặc “không giới hạn”.\\
\emph{Ví dụ:} một cổng ra của router có bộ đệm lớn, một quầy thu ngân.
\end{nhandien}

\textbf{Tham số:} $\lambda$, $\mu$. \quad \textbf{Hệ số sử dụng:} $\rho = \lambda/\mu$. \quad \textbf{Điều kiện ổn định:} $\rho<1$ (tức $\lambda<\mu$).

\bdchain{0/-/-/st, 1/\lambda/\mu/st, 2/\lambda/\mu/st, \cdots/\lambda/\mu/dots, i/\lambda/\mu/st, i{+}1/\lambda/\mu/st, \cdots/\lambda/\mu/dots}

\begin{congthuc}[Công thức M/M/1]
\renewcommand{\arraystretch}{2.1}
\begin{tabular}{@{}l l@{}}
Phân bố trạng thái & $p_i = (1-\rho)\,\rho^{\,i}$,\qquad $p_0 = 1-\rho$,\qquad $P\{\text{có} \ge k \text{ khách}\} = \rho^{\,k}$ \\
Số khách trong hệ thống & $N = \dfrac{\rho}{1-\rho} = \dfrac{\lambda}{\mu-\lambda}$ \\
Thời gian trong hệ thống & $T = \dfrac{N}{\lambda} = \dfrac{1}{\mu-\lambda} = \dfrac{1}{\mu}\cdot\dfrac{1}{1-\rho}$ \\
Thời gian chờ trong hàng & $T_Q = T - \dfrac{1}{\mu} = \dfrac{\rho}{\mu-\lambda}$ \\
Số khách trong hàng & $N_Q = \lambda T_Q = \dfrac{\rho^2}{1-\rho}$ \\
Xác suất server bận & $1-p_0 = \rho$ \\
\end{tabular}
\end{congthuc}

\begin{congthuc}[Phân bố của thời gian (với quy tắc FCFS)]
Thời gian một khách ở trong \textbf{hệ thống} có phân bố mũ với tham số $\mu-\lambda$:
\[
P\{T_{ht}\le t\} = 1-e^{-(\mu-\lambda)t}
\qquad\qquad
P\{T_{\text{chờ}} > t\} = \rho\, e^{-(\mu-\lambda)t}
\]
\end{congthuc}

\begin{vidu}[quầy tính tiền siêu thị mini (slide 33)]
Khách đến quầy theo Poisson, trung bình \textbf{40 người/giờ}. Thời gian phục vụ có phân bố mũ, trung bình \textbf{1 phút}. Có một quầy, chỗ xếp hàng đủ cho mọi người. Tìm (a) thời gian chờ trung bình; (b) xác suất từ lúc vào hàng đến lúc trả tiền xong không quá 1,5 phút.

\loigiai Đổi về phút: $\lambda = 40/60 = 2/3$ khách/phút, $\mu = 1$ khách/phút, nên $\rho = 2/3 < 1$ (ổn định).

(a) $T = \dfrac{1}{\mu-\lambda} = \dfrac{1}{1/3} = 3$ phút. Thời gian \emph{chờ trong hàng} là $T_Q = T - 1/\mu = \dapan{2 \text{ phút}}$. Nếu thầy hiểu “thời gian chờ” là $T$ thì đáp án là 3 phút.

(b) $P\{T_{ht}\le 1{,}5\} = 1-e^{-(1/3)(1{,}5)} = 1-e^{-0{,}5} \approx \dapan{0{,}393}$
\end{vidu}

% -------------------------------------------------------------
\subsection{M/M/1/K (một server, sức chứa hữu hạn)}
% -------------------------------------------------------------

\begin{nhandien}
Một server, bộ đệm \textbf{có giới hạn}: “tối đa $K$ gói”, “bộ đệm chứa $B$ gói”. Khách đến khi hệ thống đầy thì \textbf{bị từ chối/mất}.\\
Câu hỏi đặc trưng: \emph{xác suất mất gói, xác suất bị từ chối, thông lượng thực tế}.
\end{nhandien}

\textbf{Tham số:} $\lambda$, $\mu$, $K$ (sức chứa \emph{tính cả} khách đang phục vụ, bộ đệm có $K-1$ chỗ). \quad $\rho=\lambda/\mu$.\\
\textbf{Không cần điều kiện ổn định}, vì số trạng thái hữu hạn. $\rho\ge1$ vẫn tính được.

\bdchain{0/-/-/st, 1/\lambda/\mu/st, 2/\lambda/\mu/st, \cdots/\lambda/\mu/dots, K{-}1/\lambda/\mu/st, K/\lambda/\mu/st}

\begin{congthuc}[Công thức M/M/1/$K$ (với $\rho \ne 1$)]
\renewcommand{\arraystretch}{2.1}
\begin{tabular}{@{}l l@{}}
Phân bố trạng thái & $p_i = \rho^{\,i}\,p_0$,\qquad $p_0 = \dfrac{1-\rho}{1-\rho^{K+1}}$,\qquad $i=0,\dots,K$ \\
\textbf{Xác suất từ chối} & $p_K = \dfrac{(1-\rho)\,\rho^{K}}{1-\rho^{K+1}}$ \\
Tốc độ đến thực tế & $\lambda_a = \lambda\,(1-p_K)$\qquad (kiểm tra: $\lambda_a = \mu(1-p_0)$) \\
Số khách trong hệ thống & $N = \dfrac{\rho}{1-\rho} - \dfrac{(K+1)\,\rho^{K+1}}{1-\rho^{K+1}}$ \\
Số khách trong hàng & $N_Q = N - (1-p_0)$ \\
Thời gian & $T = \dfrac{N}{\lambda_a} = \dfrac{1}{\lambda}\cdot\dfrac{1}{1-p_K}\cdot N$,\qquad $T_Q = \dfrac{N_Q}{\lambda_a}$ \\
\end{tabular}
\par\medskip
\textbf{Trường hợp $\rho=1$:}\quad $p_i = \dfrac{1}{K+1}$,\quad $N = \dfrac{K}{2}$.
\end{congthuc}

\begin{luuy}[Bẫy thường gặp với M/M/1/$K$]
\begin{itemize}[leftmargin=1.2em]
  \item Dùng định lý Little thì phải dùng \textbf{$\lambda_a$}, không dùng $\lambda$. Khách bị từ chối không vào hệ thống.
  \item Đề cho “bộ đệm chứa $B$ gói” và có 1 server thì $K = B+1$.
  \item $N_Q \ne N - \rho$. Ở đây số server bận trung bình là $1-p_0$, không phải $\rho$.
\end{itemize}
\end{luuy}

\begin{vidu}[router có bộ đệm nhỏ]
Một cổng router nhận gói theo Poisson với $\lambda = 90$ gói/s. Thời gian truyền mỗi gói có phân bố mũ với $\mu=100$ gói/s. Router chứa tối đa $K=5$ gói (tính cả gói đang truyền). Tính xác suất mất gói, thông lượng, $N$, $T$, $T_Q$.

\loigiai $\rho = 0{,}9$.\quad $p_0 = \dfrac{1-0{,}9}{1-0{,}9^6} = \dfrac{0{,}1}{0{,}4686} \approx 0{,}2134$.
\[
p_5 = 0{,}9^5\,p_0 \approx \dapan{0{,}126}\ \ (\text{mất }12{,}6\%\text{ số gói}),
\qquad \lambda_a = 90(1-0{,}126) \approx 78{,}66 \text{ gói/s}
\]
\[
N = \frac{0{,}9}{0{,}1} - \frac{6\cdot0{,}9^6}{1-0{,}9^6} \approx 2{,}195 \text{ gói},\qquad
T = \frac{N}{\lambda_a} \approx \dapan{27{,}9\text{ ms}}
\]
\[
N_Q = N-(1-p_0) \approx 1{,}408 \text{ gói},\qquad T_Q = \frac{N_Q}{\lambda_a}\approx \dapan{17{,}9\text{ ms}}
\]
\emph{Kiểm tra:} $T - T_Q = 10$ ms $= 1/\mu$, khớp.
\end{vidu}

% =============================================================
\section{Các mô hình nhiều server}
% =============================================================

Với $c$ server, ở trạng thái $i$ có $\min(i,c)$ server đang bận, nên tốc độ đi là $\min(i,c)\,\mu$. Hai đại lượng cần phân biệt:
\[
a = \frac{\lambda}{\mu}\ \ (\text{cường độ tải, Erlang}) \qquad\qquad \rho = \frac{\lambda}{c\mu} = \frac{a}{c}\ \ (\text{hệ số sử dụng \emph{mỗi} server})
\]

% -------------------------------------------------------------
\subsection{M/M/c (nhiều server, hàng đợi vô hạn) và công thức Erlang C}
% -------------------------------------------------------------

\begin{nhandien}
Có \textbf{$c$ server song song} (CPU nhiều nhân, ngân hàng nhiều quầy, $c$ đường truyền), chỗ chờ không giới hạn.\\
Câu hỏi đặc trưng: \emph{xác suất phải đợi, thời gian chờ trung bình}.
\end{nhandien}

\textbf{Điều kiện ổn định:} $\rho = \dfrac{\lambda}{c\mu} < 1$.

\bdchain[15mm]{0/-/-/st, 1/\lambda/\mu/st, 2/\lambda/2\mu/st, \cdots/\lambda/3\mu/dots, c{-}1/\lambda/(c{-}1)\mu/st, c/\lambda/c\mu/st, c{+}1/\lambda/c\mu/st, \cdots/\lambda/c\mu/dots}

\begin{congthuc}[Công thức M/M/$c$]
\renewcommand{\arraystretch}{2.2}
\begin{tabular}{@{}l l@{}}
Phân bố trạng thái & $p_i = \dfrac{a^i}{i!}\,p_0\ \ (i\le c)$,\qquad $p_i = \dfrac{a^c}{c!}\,\rho^{\,i-c}\,p_0 = \dfrac{a^i}{c!\,c^{\,i-c}}\,p_0\ \ (i\ge c)$ \\
Xác suất rỗng & $p_0 = \Bigl[\displaystyle\sum_{i=0}^{c-1}\frac{a^i}{i!} + \frac{a^c}{c!\,(1-\rho)}\Bigr]^{-1}$ \\
\textbf{Erlang C} (xác suất phải đợi) & $P_Q = \displaystyle\sum_{i\ge c} p_i = \frac{a^c}{c!\,(1-\rho)}\,p_0$ \\
Số khách trong hàng & $N_Q = P_Q\,\dfrac{\rho}{1-\rho}$ \\
Thời gian chờ & $T_Q = \dfrac{N_Q}{\lambda} = \dfrac{P_Q}{c\mu-\lambda}$ \\
Thời gian trong hệ thống & $T = \dfrac{1}{\mu} + \dfrac{P_Q}{c\mu-\lambda}$ \\
Số khách trong hệ thống & $N = \lambda T = c\rho + \dfrac{\rho\,P_Q}{1-\rho} = a + N_Q$ \\
\end{tabular}
\par\smallskip
Số server bận trung bình $= a$. Phân bố thời gian chờ: $P\{T_{\text{chờ}}>t\} = P_Q\,e^{-(c\mu-\lambda)t}$.
\end{congthuc}

\begin{luuy}[Bẫy thường gặp với M/M/$c$]
Hay nhầm $\rho$ thành $\lambda/\mu$. Với $c$ server phải dùng $\rho = \lambda/(c\mu)$. Còn $\lambda/\mu = a$ có thể lớn hơn 1 mà hệ thống vẫn ổn định.\\
\emph{Tự kiểm tra:} thay $c=1$ thì $p_0 = 1-\rho$, $P_Q = \rho$, $T = 1/(\mu-\lambda)$, đúng bằng công thức M/M/1.
\end{luuy}

\begin{vidu}[CPU 4 nhân, bộ nhớ đủ lớn (slide 61)]
Tác vụ đến theo Poisson, trung bình 20 task/$\mu$s. Thời gian xử lý có phân bố mũ, trung bình 0,08 $\mu$s. CPU có 4 nhân. Tính thời gian trung bình từ lúc nhận yêu cầu đến khi hoàn thành, và số tác vụ trung bình trong RAM.

\loigiai $\lambda = 20$, $\mu = 1/0{,}08 = 12{,}5$ (task/$\mu$s), $c=4$, nên $a = 1{,}6$ và $\rho = 0{,}4<1$.

Bảng $a^i/i!$: \ 1;\ 1{,}6;\ 1{,}28;\ 0{,}6827;\ 0{,}2731 \ (với $i=0..4$). Tổng 4 số đầu là 4{,}5627.
\[
p_0 = \Bigl[4{,}5627 + \frac{0{,}2731}{0{,}6}\Bigr]^{-1} = \frac{1}{5{,}0178}\approx 0{,}1993,
\qquad P_Q = \frac{0{,}2731}{0{,}6}\cdot 0{,}1993 \approx 0{,}0907
\]
\[
N_Q = 0{,}0907\cdot\frac{0{,}4}{0{,}6}\approx 0{,}0605,\qquad
T_Q = \frac{0{,}0605}{20} \approx 0{,}0030\ \mu\text{s},\qquad
T = 0{,}08 + 0{,}0030 \approx \dapan{0{,}083\ \mu\text{s}}
\]
\[
N = \lambda T = 1{,}6 + 0{,}0605 \approx \dapan{1{,}66 \text{ tác vụ}}
\]
Đề không cho dung lượng một tác vụ nên “lượng bộ nhớ” tính bằng số tác vụ.
\end{vidu}

% -------------------------------------------------------------
\subsection{M/M/c/c (không có hàng đợi) và công thức Erlang B}
% -------------------------------------------------------------

\begin{nhandien}
$c$ kênh/server và \textbf{không có chỗ chờ}. Khách đến lúc tất cả kênh đều bận thì bị \textbf{mất ngay}.\\
\emph{Ví dụ kinh điển:} tổng đài điện thoại có $c$ trung kế. Câu hỏi đặc trưng: \emph{xác suất tắc nghẽn, cần bao nhiêu kênh}.
\end{nhandien}

\textbf{Không cần điều kiện ổn định.} Không có hàng đợi nên $N_Q = 0$, $T_Q=0$.

\bdchain[16mm]{0/-/-/st, 1/\lambda/\mu/st, 2/\lambda/2\mu/st, \cdots/\lambda/3\mu/dots, c{-}1/\lambda/(c{-}1)\mu/st, c/\lambda/c\mu/st}

\begin{congthuc}[Công thức M/M/$c$/$c$]
\renewcommand{\arraystretch}{2.2}
\begin{tabular}{@{}l l@{}}
Phân bố trạng thái & $p_i = \dfrac{a^i}{i!}\,p_0$,\qquad $p_0 = \Bigl[\displaystyle\sum_{i=0}^{c}\frac{a^i}{i!}\Bigr]^{-1}$,\qquad $i = 0,\dots,c$ \\
\textbf{Erlang B} (xác suất tắc nghẽn) & $P_B = p_c = \dfrac{a^c/c!}{\displaystyle\sum_{i=0}^{c} a^i/i!}$ \\
Tính bằng đệ quy (dễ bấm máy) & $B(0)=1,\qquad B(k) = \dfrac{a\,B(k-1)}{k + a\,B(k-1)}$,\qquad $P_B = B(c)$ \\
Lưu lượng được phục vụ & $\lambda_a = \lambda(1-P_B)$,\qquad $N = a\,(1-P_B)$ (số kênh bận trung bình) \\
Thời gian trong hệ thống & $T = 1/\mu$ \\
\end{tabular}
\end{congthuc}

\begin{vidu}[tổng đài nội bộ]
Cuộc gọi đến theo Poisson, trung bình 2 cuộc/phút. Thời gian đàm thoại trung bình 2 phút (phân bố mũ). Tổng đài có 5 kênh. (a) Tính xác suất một cuộc gọi bị chặn. (b) Cần ít nhất bao nhiêu kênh để xác suất chặn dưới 1\%?

\loigiai $a = \lambda/\mu = 2\cdot 2 = 4$ Erlang.

(a) $P_B = \dfrac{4^5/5!}{1+4+8+10{,}667+10{,}667+8{,}533} = \dfrac{8{,}533}{42{,}867}\approx \dapan{0{,}199}$. Khoảng 20\% cuộc gọi bị chặn.

(b) Tính tiếp bằng đệ quy $B(k)$:
\begin{center}
\begin{tabular}{@{}c|ccccccc@{}}
$c$ & 5 & 6 & 7 & 8 & 9 & 10 & 11\\ \hline
$P_B$ & 0,199 & 0,117 & 0,063 & 0,030 & 0,013 & \textbf{0,0053} & 0,0019
\end{tabular}
\end{center}
Vậy cần \dapan{c = 10} kênh.
\end{vidu}

% -------------------------------------------------------------
\subsection{M/M/c/K (nhiều server, sức chứa hữu hạn): mô hình tổng quát}
% -------------------------------------------------------------

\begin{nhandien}
$c$ server, hệ thống chứa tối đa $K$ khách ($c \le K$). Câu hỏi có \textbf{cả} “xác suất bị từ chối” \textbf{và} “xác suất phải đợi”.
\end{nhandien}

\bdchain[15mm]{0/-/-/st, 1/\lambda/\mu/st, \cdots/\lambda/2\mu/dots, c/\lambda/c\mu/st, c{+}1/\lambda/c\mu/st, \cdots/\lambda/c\mu/dots, K{-}1/\lambda/c\mu/st, K/\lambda/c\mu/st}

\begin{congthuc}[Công thức M/M/$c$/$K$]
\renewcommand{\arraystretch}{2.2}
\begin{tabular}{@{}l l@{}}
Phân bố trạng thái & $p_i = \dfrac{a^i}{i!}\,p_0\ \ (0\le i\le c)$,\qquad $p_i = \dfrac{a^c}{c!}\,\rho^{\,i-c}\,p_0\ \ (c\le i\le K)$ \\
Xác suất rỗng & $p_0 = \Bigl[\displaystyle\sum_{i=0}^{c-1}\frac{a^i}{i!} + \frac{a^c}{c!}\sum_{j=0}^{K-c}\rho^{\,j}\Bigr]^{-1}$,\qquad $\displaystyle\sum_{j=0}^{K-c}\rho^{\,j} = \frac{1-\rho^{K-c+1}}{1-\rho}$ \\
\textbf{Xác suất từ chối} & $p_K = \dfrac{a^c}{c!}\,\rho^{\,K-c}\,p_0$ \\
\textbf{Xác suất phải đợi} & $P_{\text{đợi}} = \displaystyle\sum_{i=c}^{K-1} p_i = \frac{a^c}{c!}\cdot\frac{1-\rho^{K-c}}{1-\rho}\cdot p_0$ \\
Số khách trong hàng & $N_Q = \displaystyle\sum_{i=c}^{K}(i-c)\,p_i$ \\
Các đại lượng còn lại & $\lambda_a = \lambda(1-p_K)$,\quad $N = N_Q + \dfrac{\lambda_a}{\mu}$,\quad $T = \dfrac{N}{\lambda_a}$,\quad $T_Q = \dfrac{N_Q}{\lambda_a}$ \\
\end{tabular}
\par\smallskip
Khi $\rho = 1$ thì thay $\sum_{j=0}^{K-c}\rho^j$ bằng $K-c+1$. \textbf{Không cần điều kiện ổn định.}
\end{congthuc}

\begin{vidu}[CPU 4 nhân, cache tối đa 10 tác vụ (slide 61)]
Cùng số liệu với ví dụ M/M/4 ở trên nhưng hệ thống chỉ chứa tối đa $K=10$ tác vụ. Tính xác suất một tác vụ bị từ chối và xác suất phải đợi.

\loigiai $a=1{,}6$, $\rho=0{,}4$, $c=4$, $K=10$, nên $K-c = 6$.
\[
p_0 = \Bigl[4{,}5627 + 0{,}2731\cdot\frac{1-0{,}4^7}{0{,}6}\Bigr]^{-1} = \frac{1}{5{,}0170}\approx 0{,}1993
\]
\[
p_{10} = 0{,}2731\cdot 0{,}4^6\cdot 0{,}1993 \approx \dapan{2{,}23\times10^{-4}}
\qquad
P_{\text{đợi}} = 0{,}2731\cdot\frac{1-0{,}4^6}{0{,}6}\cdot0{,}1993 \approx \dapan{0{,}0903}
\]
Nếu “10 tác vụ” chỉ là phần chờ trong cache thì $K = 10+4 = 14$. Hãy đọc kỹ đề.
\end{vidu}

% -------------------------------------------------------------
\subsection{\texorpdfstring{M/M/$\infty$}{M/M/∞} (vô số server)}
% -------------------------------------------------------------

Mỗi khách đến đều được phục vụ ngay, không ai phải chờ.\quad $\mu_i = i\mu$.
\begin{congthuc}[Công thức M/M/$\infty$]
\[
p_i = \frac{a^i}{i!}\,e^{-a}\ \ (\text{phân bố Poisson}),\qquad N = a = \frac{\lambda}{\mu},\qquad T = \frac{1}{\mu},\qquad N_Q = T_Q = 0
\]
\end{congthuc}
\emph{Ứng dụng:} mô hình hóa số người dùng đang online, số phiên đang mở, khi tài nguyên coi như không giới hạn.

% =============================================================
\section{Mô hình mở rộng và quy tắc phục vụ}
% =============================================================

\begin{luuy}[Phạm vi]
Mục~\ref{sec:mg1} (M/G/1, M/D/1) \textbf{không có trong slide} của môn. Phần này để tham khảo và để hiểu vì sao giả thiết “phục vụ mũ” lại quan trọng. Mục~\ref{sec:uutien} tóm tắt phần lý thuyết về quy tắc phục vụ có trong slide.
\end{luuy}

\subsection{M/G/1 và M/D/1}\label{sec:mg1}

Khi thời gian phục vụ $X$ \textbf{không} có phân bố mũ thì không còn chuỗi sinh--tử. Lúc đó dùng công thức \textbf{Pollaczek--Khinchine (P-K)}:

\begin{congthuc}[{Công thức P-K cho M/G/1 ($\rho = \lambda E[X] < 1$)}]
\[
T_Q = \frac{\lambda\,E[X^2]}{2(1-\rho)},\qquad E[X^2] = \operatorname{Var}(X) + \bigl(E[X]\bigr)^2,\qquad
T = T_Q + E[X],\qquad N_Q = \lambda T_Q
\]
\end{congthuc}

\begin{center}
\begin{tabular}{@{}l c c@{}}
\toprule
\textbf{Mô hình} & \textbf{$E[X^2]$} & \textbf{$T_Q$} \\
\midrule
M/M/1 (mũ) & $2/\mu^2$ & $\dfrac{\rho}{\mu(1-\rho)}$ \\[2mm]
M/D/1 (hằng số, ví dụ gói cùng độ dài) & $1/\mu^2$ & $\dfrac{\rho}{2\mu(1-\rho)}$ \\
\bottomrule
\end{tabular}
\end{center}

\textbf{Nhận xét:} với cùng $\rho$, thời gian chờ của \textbf{M/D/1 chỉ bằng một nửa M/M/1}. Thời gian phục vụ càng biến động thì hàng đợi càng dài. Ví dụ với $\rho=0{,}8$ và $\mu = 1$: M/M/1 có $T_Q = 4$, M/D/1 có $T_Q = 2$.

\subsection{Quy tắc phục vụ và hàng đợi ưu tiên}\label{sec:uutien}

\begin{center}
\begin{tabularx}{\linewidth}{@{}l X@{}}
\toprule
\textbf{Quy tắc} & \textbf{Mô tả} \\
\midrule
FCFS / FIFO & Đến trước phục vụ trước. Đây là mặc định khi đề không nói gì \\
LCFS / LIFO & Đến sau phục vụ trước (kiểu ngăn xếp) \\
SIRO (RANDOM) & Chọn ngẫu nhiên, mọi khách đang chờ có xác suất được chọn như nhau \\
SJF / SPF & Ưu tiên công việc có thời gian xử lý ngắn nhất \\
RR (Round Robin) & Mỗi khách được phục vụ một lát thời gian cố định, chưa xong thì quay lại cuối hàng \\
PS (Processor Sharing) & Mọi khách chia đều năng lực phục vụ \\
FB (Foreground--Background) & Cố gắng đạt hiệu quả như SJF mà không cần biết trước thời gian phục vụ \\
\bottomrule
\end{tabularx}
\end{center}

\textbf{Thuật toán xếp hàng trong router:}
\begin{itemize}
  \item \textbf{FIFO Queuing}: đơn giản, không cần cấu hình, nhưng không phân biệt ưu tiên. Một nguồn phát quá nhiều có thể làm tăng trễ của mọi luồng khác.
  \item \textbf{PQ (Priority Queuing)}: nhiều hàng có mức ưu tiên khác nhau, hàng ưu tiên cao luôn được phục vụ trước. Không tự thích nghi khi điều kiện mạng thay đổi.
  \item \textbf{CQ (Custom Queuing)}: mỗi loại lưu lượng có một phần hàng đợi riêng, phục vụ xoay vòng (round-robin) để chia băng thông theo tỷ lệ. Không tự thích nghi.
  \item \textbf{WFQ (Weighted Fair Queuing)}: giống CQ nhưng trọng số được hệ thống tự điều chỉnh (qua báo hiệu QoS), tự thích ứng với lưu lượng.
\end{itemize}

\textbf{Hàng đợi có ưu tiên}, chia theo cách xử lý khi khách ưu tiên cao đến:
\begin{itemize}
  \item \textbf{Không ưu tiên phục vụ trước} (non-preemptive, HOL): khách ưu tiên cao vẫn phải đợi khách đang được phục vụ xong.
  \item \textbf{Ưu tiên phục vụ trước} (preemptive): ngắt ngay khách ưu tiên thấp. Có ba kiểu: \emph{phục hồi} (resume, làm tiếp từ chỗ dừng), \emph{không lấy mẫu lại} (làm lại từ đầu với thời gian cũ), \emph{lấy mẫu lại} (làm lại với thời gian phục vụ mới).
\end{itemize}

% =============================================================
\section{Mạng hàng đợi}
% =============================================================

Mạng gồm nhiều nút hàng đợi nối với nhau. Sau khi được phục vụ ở một nút, khách chuyển sang nút khác hoặc rời mạng theo các xác suất định tuyến.

\begin{center}
\begin{tabularx}{\linewidth}{@{}c X@{}}
\toprule
\textbf{Ký hiệu} & \textbf{Ý nghĩa} \\
\midrule
$\lambda$ & Tổng tốc độ đến trung bình từ bên ngoài vào mạng \\
$\mu_i$ & Tốc độ phục vụ của nút $i$ \\
$r_{sj}$ & Xác suất khách từ nguồn được đưa vào nút $j$ \\
$r_{jk}$ & Xác suất khách rời nút $j$ đi tiếp tới nút $k$ \\
$r_{jd}$ & Xác suất khách rời nút $j$ đi ra đích (rời mạng) \\
$\theta_i$ & \textbf{Thông lượng} trung bình qua nút $i$ (tốc độ đến tổng tại nút $i$) \\
\bottomrule
\end{tabularx}
\end{center}

\subsection{Phương trình lưu lượng}

\begin{congthuc}[Cân bằng luồng tại mỗi nút]
\[
\theta_j = \lambda\, r_{sj} + \sum_{k} \theta_k\, r_{kj}\qquad\text{với mọi nút } j
\]
Luồng vào nút $j$ = luồng đến từ ngoài + tổng các luồng chuyển tới từ những nút khác (kể cả từ chính nó).\\
\textbf{Điều kiện ổn định của mạng:} $\rho_i = \theta_i/\mu_i < 1$, tức $\mu_i > \theta_i$, với \textbf{mọi} nút $i$.
\end{congthuc}

\subsection{Định lý Burke}

Đầu ra của một hàng M/M/1 (hoặc M/M/$c$) \textbf{ổn định} có tốc độ $\lambda$ cũng là một \textbf{luồng Poisson với tốc độ $\lambda$}. Hơn nữa, trạng thái của hàng đợi tại thời điểm $t$ độc lập với quá trình rời đi trước $t$.\\
\emph{Hệ quả:} với mạng nối tiếp (không có vòng phản hồi), đầu vào của mọi nút đều là Poisson. Nhờ vậy có thể phân tích từng nút như một M/M/1 riêng.

\subsection{Mạng Jackson và định lý Jackson}

\textbf{Mạng Jackson} (mạng mở) có các đặc điểm: mỗi nút là hàng FIFO, hàng đợi vô hạn, thời gian phục vụ có phân bố mũ $\operatorname{Exp}(\mu_i)$, khách từ ngoài đến theo Poisson, định tuyến theo xác suất cố định.

\begin{congthuc}[Định lý Jackson: coi mỗi nút là một M/M/1 độc lập]
\[
p(n_1,n_2,\dots,n_M) = \prod_{i=1}^{M} (1-\rho_i)\,\rho_i^{\,n_i},\qquad \rho_i = \frac{\theta_i}{\mu_i}
\]
\[
E(N_i) = \frac{\rho_i}{1-\rho_i},\qquad E(T_i) = \frac{1}{\theta_i}\cdot\frac{\rho_i}{1-\rho_i} = \frac{1}{\mu_i-\theta_i}
\]
\[
N_{\text{mạng}} = \sum_i E(N_i),\qquad T_{\text{mạng}} = \frac{N_{\text{mạng}}}{\lambda} = \sum_i \frac{\theta_i}{\lambda}\,E(T_i),\qquad P\{\text{mạng rỗng}\} = \prod_i (1-\rho_i)
\]
\end{congthuc}
$\theta_i/\lambda$ là \textbf{số lần trung bình một khách đi qua nút $i$}. Nút có vòng phản hồi thì giá trị này lớn hơn 1.

\begin{luuy}[Bẫy khi tính trễ toàn mạng]
Trễ toàn mạng là $N_{\text{mạng}}$ chia cho \textbf{$\lambda$} (tốc độ đến từ \emph{bên ngoài}), \textbf{không} chia cho $\sum\theta_i$. Có vòng phản hồi thì một khách đi qua một nút nhiều lần, nên $\sum \theta_i > \lambda$.
\end{luuy}

\begin{vidu}[mạng hai nút có phản hồi]
Gói đến nút 1 theo Poisson với $\lambda = 10$ gói/s. Xong nút 1 thì sang nút 2. Xong nút 2 thì gói quay lại nút 1 với xác suất 0,5, hoặc rời mạng với xác suất 0,5. Cho $\mu_1 = 25$, $\mu_2 = 50$ gói/s.

\begin{center}
\begin{tikzpicture}[
  nd/.style={draw=accent,very thick,fill=accentsoft,rounded corners=2pt,minimum width=24mm,minimum height=10mm,align=center,font=\small},
  ar/.style={-{Stealth[length=2.2mm]},thick,color=ink!80}]
  \node (in) at (0,0) {};
  \node[nd] (q1) at (3,0) {Nút 1\\[-2pt]\footnotesize$\mu_1=25$};
  \node[nd] (q2) at (7,0) {Nút 2\\[-2pt]\footnotesize$\mu_2=50$};
  \node (out) at (10.5,0) {};
  \draw[ar] (in) -- node[above,font=\small]{$\lambda=10$} (q1);
  \draw[ar] (q1) -- node[above,font=\small]{$\theta_1$} (q2);
  \draw[ar] (q2) -- node[above,font=\small]{0,5 ra} (out);
  \draw[ar] (q2.south) -- ++(0,-0.7) -| node[pos=0.25,below,font=\small]{0,5 quay lại} (q1.south);
\end{tikzpicture}
\end{center}

\loigiai Phương trình lưu lượng: $\theta_1 = 10 + 0{,}5\,\theta_2$, $\theta_2 = \theta_1$. Suy ra $\theta_1=\theta_2 = 20$ gói/s.\\
$\rho_1 = 20/25 = 0{,}8$,\quad $\rho_2 = 20/50 = 0{,}4$. Cả hai đều nhỏ hơn 1 nên mạng ổn định.
\[
E(N_1) = \frac{0{,}8}{0{,}2} = 4,\qquad E(N_2) = \frac{0{,}4}{0{,}6}\approx 0{,}667,\qquad N_{\text{mạng}}\approx 4{,}667 \text{ gói}
\]
\[
T_{\text{mạng}} = \frac{4{,}667}{10} \approx \dapan{0{,}467\text{ s}},\qquad P\{\text{mạng rỗng}\} = 0{,}2\cdot 0{,}6 = \dapan{0{,}12}
\]
\emph{Kiểm tra:} mỗi gói qua mỗi nút trung bình $20/10 = 2$ lần, nên $T = 2\cdot\frac{1}{25-20} + 2\cdot\frac{1}{50-20} = 0{,}4 + 0{,}067 = 0{,}467$ s. Khớp.
\end{vidu}

\subsection{Xấp xỉ độc lập Kleinrock (KIA)}

Dùng để tính trễ trong \textbf{mạng truyền dữ liệu nhiều luồng}. Mỗi liên kết $(i,j)$ được coi là một hàng M/M/1 độc lập, dù trên thực tế độ dài gói giữ nguyên khi đi qua các liên kết.

\begin{center}
\begin{tabularx}{\linewidth}{@{}c X@{}}
\toprule
$x_s$ & Tốc độ gói của luồng $s$ (đi theo một đường cố định từ nguồn tới đích) \\
$\lambda_{ij}$ & Tốc độ đến trên liên kết $(i,j)$ = tổng $x_s$ của các luồng đi qua liên kết đó. Nếu luồng $s$ chỉ đưa một phần $f_{ij}(s)$ qua liên kết thì $\lambda_{ij} = \sum_s f_{ij}(s)\,x_s$ \\
$\mu_{ij}$ & Tốc độ phục vụ của liên kết: $\mu_{ij} = C_{ij}/L$ (dung lượng chia cho độ dài gói trung bình) \\
$\gamma$ & Tổng tốc độ đến của toàn mạng: $\gamma = \sum_s x_s$ \\
$d_{ij}$ & Trễ lan truyền trên liên kết $(i,j)$ \\
\bottomrule
\end{tabularx}
\end{center}

\begin{congthuc}[Công thức KIA]
\[
N_{ij} = \frac{\lambda_{ij}}{\mu_{ij}-\lambda_{ij}},\qquad
N = \sum_{(i,j)} \frac{\lambda_{ij}}{\mu_{ij}-\lambda_{ij}},\qquad
T = \frac{1}{\gamma}\sum_{(i,j)}\frac{\lambda_{ij}}{\mu_{ij}-\lambda_{ij}} = \sum_{(i,j)}\alpha_{ij}T_{ij}
\]
với $\alpha_{ij} = \lambda_{ij}/\gamma$ và $T_{ij} = \dfrac{1}{\mu_{ij}-\lambda_{ij}}$. Nếu tính thêm \textbf{trễ lan truyền}:
\[
T = \frac{1}{\gamma}\sum_{(i,j)}\Bigl(\frac{\lambda_{ij}}{\mu_{ij}-\lambda_{ij}} + \lambda_{ij}\,d_{ij}\Bigr)
\]
\textbf{Trễ của một luồng} đi theo đường $p$:\quad
$T_p = \displaystyle\sum_{(i,j)\in p}\Bigl(\frac{\rho_{ij}}{\mu_{ij}-\lambda_{ij}} + \frac{1}{\mu_{ij}} + d_{ij}\Bigr)$, với $\rho_{ij} = \lambda_{ij}/\mu_{ij}$.
\end{congthuc}

\begin{vidu}[ba nút A--B--C]
Liên kết A$\to$B có $\mu=500$ gói/s, liên kết B$\to$C có $\mu = 400$ gói/s. Có ba luồng: $s_1$: A$\to$C với 100 gói/s; $s_2$: A$\to$B với 150 gói/s; $s_3$: B$\to$C với 200 gói/s. Bỏ qua trễ lan truyền.

\loigiai $\lambda_{AB} = 100+150 = 250$,\quad $\lambda_{BC} = 100 + 200 = 300$,\quad $\gamma = 450$ gói/s.
\[
N_{AB} = \frac{250}{500-250} = 1,\qquad N_{BC} = \frac{300}{400-300} = 3,\qquad T = \frac{1+3}{450}\approx\dapan{8{,}89\text{ ms}}
\]
Trễ riêng của luồng $s_1$ (đi qua cả hai liên kết): $T_{s_1} = \dfrac{1}{250} + \dfrac{1}{100} = 4 + 10 = \dapan{14\text{ ms}}$.
\end{vidu}

% =============================================================
\section{Tổng kết và ôn thi}
% =============================================================

\subsection{Bảng so sánh tất cả các mô hình}

\begingroup
\small
\everymath{\displaystyle}
\renewcommand{\arraystretch}{2.6}
\begin{tabularx}{\linewidth}{@{}>{\bfseries}l >{\raggedright\arraybackslash}p{2.3cm} >{\raggedright\arraybackslash}X >{\raggedright\arraybackslash}X@{}}
\toprule
\normalfont\textbf{Mô hình} & \textbf{$\rho$ / điều kiện} & \textbf{$p_0$} & \textbf{Kết quả đặc trưng} \\
\midrule
M/M/1 & $\rho = \frac{\lambda}{\mu}<1$ & $1-\rho$ & $N = \frac{\rho}{1-\rho}$,\ $T = \frac{1}{\mu-\lambda}$,\ $T_Q = \frac{\rho}{\mu-\lambda}$,\ $N_Q = \frac{\rho^2}{1-\rho}$ \\
M/M/1/$K$ & $\rho=\frac{\lambda}{\mu}$, không cần $<1$ & $\frac{1-\rho}{1-\rho^{K+1}}$ & $p_K = \rho^K p_0$;\ $N = \frac{\rho}{1-\rho}-\frac{(K+1)\rho^{K+1}}{1-\rho^{K+1}}$;\ $T = \frac{N}{\lambda_a}$ \\
M/M/$c$ & $\rho = \frac{\lambda}{c\mu}<1$ & $\Bigl[\sum_{i=0}^{c-1}\frac{a^i}{i!}+\frac{a^c}{c!(1-\rho)}\Bigr]^{-1}$ & $P_Q = \frac{a^c p_0}{c!(1-\rho)}$;\ $T_Q = \frac{P_Q}{c\mu-\lambda}$;\ $N = a + N_Q$ \\
M/M/$c$/$c$ & $a = \frac{\lambda}{\mu}$, không cần điều kiện & $\Bigl[\sum_{i=0}^{c}\frac{a^i}{i!}\Bigr]^{-1}$ & $P_B = \frac{a^c}{c!}p_0$ (Erlang B);\ $T = \frac1\mu$;\ $N = a(1-P_B)$ \\
M/M/$c$/$K$ & $\rho = \frac{\lambda}{c\mu}$, không cần $<1$ & $\Bigl[\sum_{i=0}^{c-1}\frac{a^i}{i!}+\frac{a^c}{c!}\sum_{j=0}^{K-c}\rho^j\Bigr]^{-1}$ & $p_K = \frac{a^c}{c!}\rho^{K-c}p_0$;\ $P_{\text{đợi}}=\sum_{i=c}^{K-1}p_i$;\ dùng $\lambda_a$ \\
M/M/$\infty$ & không cần điều kiện & $e^{-a}$ & $p_i = \frac{a^i}{i!}e^{-a}$;\ $N = a$;\ $T = \frac1\mu$ \\
M/G/1 & $\rho = \lambda E[X]<1$ & $1-\rho$ & $T_Q = \frac{\lambda E[X^2]}{2(1-\rho)}$ (P-K) \\
Jackson & $\frac{\theta_i}{\mu_i}<1\ \forall i$ & $\prod_i(1-\rho_i)$ & $N = \sum\frac{\rho_i}{1-\rho_i}$;\ $T = \frac{N}{\lambda}$ \\
KIA & $\lambda_{ij}<\mu_{ij}$ & -- & $T = \frac1\gamma\sum\frac{\lambda_{ij}}{\mu_{ij}-\lambda_{ij}}$ \\
\bottomrule
\end{tabularx}
\endgroup

\subsection{Nhận diện mô hình từ đề bài}\label{sec:nhandien}

\begin{center}
\begin{tabularx}{\linewidth}{@{}X l@{}}
\toprule
\textbf{Dấu hiệu trong đề} & \textbf{Kết luận} \\
\midrule
“đến theo tiến trình Poisson”, “khoảng cách giữa hai lần đến có phân bố mũ” & A = M \\
“thời gian phục vụ/xử lý có phân bố mũ”, “độ dài gói có phân bố mũ” & S = M \\
“thời gian phục vụ cố định”, “gói có độ dài cố định” & S = D \\
“một quầy”, “một đường truyền”, “một CPU” & $c=1$ \\
“4 nhân”, “3 quầy”, “$c$ kênh” & $c = 4,\ 3,\ \dots$ \\
“đủ chỗ cho tất cả”, “bộ nhớ đủ lớn”, “không giới hạn” & $K=\infty$ \\
“tối đa $K$ khách trong hệ thống” & $K$ \\
“bộ đệm (chỗ chờ) chứa $B$” & $K = B + c$ \\
“không có chỗ chờ”, “bận thì cuộc gọi bị mất” & $K = c$ (Erlang B) \\
Hỏi “xác suất phải đợi” với nhiều server & Erlang C (M/M/$c$) \\
Hỏi “xác suất từ chối/mất gói/tắc nghẽn” & $K$ hữu hạn, tính $p_K$ \\
Nhiều router/nút nối nhau, có xác suất định tuyến & Mạng Jackson \\
Nhiều luồng đi qua các liên kết có dung lượng $C_{ij}$ & KIA \\
\bottomrule
\end{tabularx}
\end{center}

\subsection{Mười lỗi hay mất điểm}

\begin{luuy}[Kiểm tra trước khi nộp bài]
\begin{enumerate}[leftmargin=1.5em]
  \item \textbf{Không đổi đơn vị}: $\lambda$ theo giờ nhưng thời gian phục vụ theo phút; Byte và bit ($1$ Byte $= 8$ bit); ms và s.
  \item Đề cho dung lượng $C$ (bps) và độ dài gói $L$ (bit) mà quên tính \textbf{$\mu = C/L$}.
  \item “Trung bình 5 phút/chuyến” là \textbf{$1/\lambda$}, không phải $\lambda$.
  \item Với $c$ server dùng nhầm $\rho = \lambda/\mu$. Đúng phải là \textbf{$\rho = \lambda/(c\mu)$}.
  \item Hệ thống có $K$ hữu hạn mà vẫn dùng $\lambda$ trong định lý Little. Phải dùng \textbf{$\lambda_a = \lambda(1-p_K)$}.
  \item Nhầm \textbf{$T$} (trong hệ thống) với \textbf{$T_Q$} (trong hàng đợi).
  \item Coi “bộ đệm $B$” là $K$, trong khi $K = B + c$.
  \item Nhầm \textbf{Erlang B} (không có hàng đợi, xác suất bị chặn) với \textbf{Erlang C} (có hàng đợi, xác suất phải chờ).
  \item Mạng có phản hồi: tính trễ toàn mạng phải chia cho $\lambda$ \textbf{từ bên ngoài}, không chia cho $\sum\theta_i$.
  \item Quên \textbf{kiểm tra điều kiện ổn định} trước khi dùng công thức của hệ thống vô hạn.
\end{enumerate}
\end{luuy}

\subsection{Mẹo tính nhanh khi đi thi}

\begin{itemize}
  \item Lập sẵn \textbf{bảng $a^i/i!$} cho $i = 0,\dots,c$. Mọi công thức M/M/$c$, M/M/$c$/$c$, M/M/$c$/$K$ đều dùng lại bảng này.
  \item Erlang B: dùng đệ quy $B(k) = \dfrac{aB(k-1)}{k+aB(k-1)}$ thay vì tính giai thừa lớn.
  \item Tổng cấp số nhân: $\displaystyle\sum_{j=0}^{n}\rho^j = \frac{1-\rho^{n+1}}{1-\rho}$.
  \item \textbf{Tự kiểm tra} bằng ba hệ thức: $T - T_Q = 1/\mu$;\ $N - N_Q$ = số server bận trung bình ($= a$ khi không từ chối, $= \lambda_a/\mu$ khi có từ chối);\ $\sum p_i = 1$.
\end{itemize}


\end{document}
